\section{Nonlinear coordinate changes of Stieltjes finite parts}
\label{nc:section}

The incoming Mellin--dilation and Stieltjes--harmonic reports explicitly
ask for the replacement of a local affine coordinate $qt$ by
$f(t)=qt+c_2t^2+\cdots$, including every Stieltjes index, every argument
derivative, and a composition law \cite{incoming-md,incoming-sh}.
The answer is a finite residue formula. Besides recovering their affine
polynomial, it yields a sharp additional phenomenon: if $f'(0)=1$, the
finite part of $\gamma_n^{(r)}$ transforms without a local correction
whenever $n\geq2r$.

The method below is direct analytic regularization. Its general
finite-part principle is not presented as a claim of historical priority.
The contribution to the stated source question is the complete formula,
its specialization to the Stieltjes hierarchy, the exact composition law,
and the sharp filtration of the surviving local terms.

\subsection{A universal residue formula}

All coordinates in this section are local real coordinates, with the
singular approach from the right. Let
\begin{equation}\label{nc:loggerm}
 h(x)=\sum_{\ell=0}^{L}a_\ell(x)(\log x)^\ell,\qquad x>0,
\end{equation}
where the $a_\ell$ are meromorphic germs at zero, each with a finite pole
order. Define $\FP_xh$ by the constant term of
$\int_\eps h(x)\phi(x)\dd x$ as $\eps\downarrow0$, for a smooth
compactly supported test function $\phi$. Every negative power of
$\eps$ and every nonconstant power of $\log\eps$ is discarded.
The positive-side germ can be extended by zero on the left; a fixed
smooth completion can equally well be added. Such a completion changes
neither the residue nor the coordinate correction.

Let $f$ be a real-analytic local diffeomorphism with
\begin{equation}\label{nc:coordinate}
 f(0)=0,\qquad f'(0)=q>0,
 \qquad \lambda_f(t)=\log\frac{f(t)}t.
\end{equation}
The last function is analytic at zero, with
$\lambda_f(0)=\log q$. Pullback here means scalar distributional
pullback:
\begin{equation}\label{nc:pullback}
 \langle f^*T,\phi\rangle
 =\left\langle T,
       (\phi\circ f^{-1})(f^{-1})'\right\rangle.
\end{equation}
In particular, the pullback of an ordinary function is its composition
with $f$.

\begin{theorem}[Nonlinear finite-part coordinate law]
\label{nc:universal}
For the germ \eqref{nc:loggerm}, define
\begin{equation}\label{nc:residue-anomaly}
 \left\langle\mathfrak A_f[h],\phi\right\rangle
 =\sum_{\ell=0}^{L}\frac1{\ell+1}
  \Res_{t=0}
    a_\ell(f(t))\lambda_f(t)^{\ell+1}\phi(t)\dd t.
\end{equation}
Then
\begin{equation}\label{nc:universal-law}
 \boxed{\quad
 f^*\FP_xh=\FP_t(h\circ f)+\mathfrak A_f[h].
 \quad}
\end{equation}
If every $a_\ell$ has pole order at most $R$, the correction is a
linear combination of
$\delta_0,\delta_0',\ldots,\delta_0^{(R-1)}$; for $R=0$ it is zero.
Every coefficient is determined by finite Taylor data. The residue
notation for a smooth $\phi$ means its finite Taylor coefficient;
holomorphic extension of the test function is unnecessary.
\end{theorem}

\begin{proof}
It suffices to take $h(x)=a(x)\log^\ell x$. For $\Re z$ sufficiently
large, the pairing of $x_+^za(x)$ with a test function is an ordinary
integral. If $a(x)=x^{-R}A(x)$ with $A$ holomorphic, subtract the Taylor
polynomial of $A(x)\phi(x)$ through degree $R-1$. The remainder is
locally integrable for $z$ near zero, uniformly on a small closed disk.
This constructs a meromorphic distribution with at most a simple pole
at zero, and
\begin{equation}\label{nc:regulator-fp}
 \FP_x\bigl(a(x)\log^\ell x\bigr)
       =\ell![z^\ell]\bigl(x_+^za(x)\bigr).
\end{equation}
Here the bracket is a Laurent coefficient. To check the normalization,
the only exceptional Taylor monomial is $x^{z-1}$, whose continued
integral to a fixed positive $h_0$ is $h_0^z/z$. Its coefficient
$[z^\ell]$ is $\log^{\ell+1}h_0/(\ell+1)!$, precisely the constant
remaining after subtracting the lower-endpoint logarithm. Every other
Taylor monomial has a denominator nonzero at $z=0$.

After pullback the regulated integrand is
\[
 t^z a(f(t))e^{z\lambda_f(t)}\phi(t).
\]
Its polar Taylor term at zero has numerator
\begin{equation}\label{nc:Bregulator}
 B(z)=\Res_{t=0}a(f(t))e^{z\lambda_f(t)}\phi(t)\dd t.
\end{equation}
Thus the exceptional contribution is $B(z)h_0^z/z$. On the other
hand, the cutoff integral of the same term is
$B(z)(h_0^z-\eps^z)/z$. In its coefficient $[z^\ell]$, the lower
endpoint has constant term $[z^{\ell+1}]B(z)$: all its other terms
are positive powers of $\log\eps$. Consequently the meromorphic
coefficient exceeds the coefficientwise cutoff constant by
\[
 \ell![z^{\ell+1}]B(z)
   =\frac1{\ell+1}\Res_{t=0}
      a(f(t))\lambda_f(t)^{\ell+1}\phi(t)\dd t.
\]
Nonpolar Taylor terms have lower-endpoint powers different from zero
and introduce no constant discrepancy. The subtracted remainder has
uniform integrable majorants, so Cauchy's coefficient formula justifies
the coefficient extraction there. This proves
\eqref{nc:universal-law}.

Finally, $a_\ell(f(t))$ has pole order at most $R$ and $\lambda_f$ is
analytic. Only derivatives $\phi^{(j)}(0)$ with $j\leq R-1$ can occur
in its residue. This proves both the support and finite-data claims.
\end{proof}

\subsection{Composition and the transformed singular germ}

\begin{theorem}[Exact residue cocycle]
\label{nc:composition}
For two orientation-preserving analytic coordinate germs $f,g$,
\begin{equation}\label{nc:cocycle}
 \boxed{\quad
 \mathfrak A_{g\circ f}[h]
    =f^*\mathfrak A_g[h]+\mathfrak A_f[h\circ g].
 \quad}
\end{equation}
The last term acts on the transformed germ $h\circ g$, which again has
the form \eqref{nc:loggerm}. This dependence is part of the composition
law.
\end{theorem}

\begin{proof}
The assertion follows by applying Theorem~\ref{nc:universal} twice.
There is also an explicit coefficient proof. For a single term
$h(x)=a(x)\log^\ell x$, put
$A=\lambda_f(t)$ and $B=\lambda_g(f(t))$. Then
\[
 \lambda_{g\circ f}=A+B,\qquad
 h(g(y))=\sum_{k=0}^\ell\binom\ell k
   a(g(y))\lambda_g(y)^{\ell-k}(\log y)^k.
\]
Residue substitution $y=f(t)$ in the pullback
\eqref{nc:pullback} shows that $f^*\mathfrak A_g[h]$ contributes the
factor $B^{\ell+1}/(\ell+1)$. The other term contributes
\[
 \sum_{k=0}^\ell\frac{\binom\ell k}{k+1}B^{\ell-k}A^{k+1}
 =\frac{(A+B)^{\ell+1}-B^{\ell+1}}{\ell+1}.
\]
The sum is precisely \eqref{nc:residue-anomaly} for $g\circ f$.
This finite polynomial identity proves composition with all scale and
nonlinear terms retained.
\end{proof}

For later implementation, the same residue convention gives
\begin{equation}\label{nc:delta-pullback}
 \langle f^*\delta_0^{(r)},\phi\rangle
   =(-1)^rr!\Res_{t=0}f(t)^{-r-1}\phi(t)\dd t.
\end{equation}
Indeed, substitution $x=f(t)$ changes the residue into the coefficient
of $x^r$ in $(\phi\circ f^{-1})(f^{-1})'$, in agreement with the
definition of $\delta_0^{(r)}$. Equations
\eqref{nc:cocycle} and \eqref{nc:delta-pullback} provide an exact
composition algorithm. A formula that retained $h$ unchanged in the
last term of \eqref{nc:cocycle} would generally be incorrect.

\subsection{Every Stieltjes index and argument derivative}

The local singularity is fixed by the Hurwitz convention:
\begin{equation}\label{nc:Stieltjes-local}
 \gamma_n(x)=\frac{\log^n x}{x}+\gamma_n(1+x),
 \qquad \gamma_0(x)=-\psi(x).
\end{equation}
The second term is analytic near zero. Set
\begin{equation}\label{nc:Tr}
 T_{n,r}=\FP_x\gamma_n^{(r)}(x),\qquad G_n=T_{n,0}.
\end{equation}
Locally these may be regarded as the periodic finite parts used in
\cite{incoming-md,incoming-sh}; the smooth completion in
\eqref{nc:Stieltjes-local} is then the same on both sides of the
endpoint. Every comparison involving $D^rG_n$ below uses this common
smooth periodic completion. In particular, the analytic remainder is
smooth across zero and contributes no step-function derivatives.

For $r\geq0$, let
\begin{align}
 \mathcal E_r(z)&=\prod_{j=1}^r(1-z/j),\qquad \mathcal E_0(z)=1,
       \label{nc:Er}\\
 b_{n,r}(L)&=n![z^{n+1}]e^{Lz}\mathcal E_r(z),\notag\\
 \alpha_{n,r}(L)&=b_{n,r}(L)-b_{n,r}(0)
       =n![z^{n+1}](e^{Lz}-1)\mathcal E_r(z).
       \label{nc:alpha}
\end{align}
If $e_j(r)$ is the elementary symmetric polynomial of degree $j$ in
$1,1/2,\ldots,1/r$, with $e_0(r)=1$ and $e_j(r)=0$ for $j>r$, then
\begin{equation}\label{nc:alpha-finite}
 \alpha_{n,r}(L)=n!\sum_{j=0}^{\min(r,n)}
      \frac{(-1)^je_j(r)}{(n+1-j)!}L^{n+1-j}.
\end{equation}
The subtraction of $b_{n,r}(0)$ is necessary because $T_{n,r}$ in
\eqref{nc:Tr} is the finite part of a pointwise derivative.

\begin{theorem}[Nonlinear Stieltjes transport]
\label{nc:Stieltjes}
For every $n,r\geq0$,
\begin{equation}\label{nc:Stieltjes-law}
 f^*T_{n,r}=\FP_t\gamma_n^{(r)}(f(t))+\mathfrak A_{n,r}(f),
\end{equation}
where
\begin{equation}\label{nc:Stieltjes-residue}
 \langle\mathfrak A_{n,r}(f),\phi\rangle
   =(-1)^rr!\Res_{t=0}
     f(t)^{-r-1}\alpha_{n,r}(\lambda_f(t))\phi(t)\dd t.
\end{equation}
In particular,
\begin{align}
 \mathfrak A_{n,r}(f)&=\sum_{j=0}^r c_{n,r,j}(f)\delta_0^{(j)},
       \label{nc:finite-delta}\\
 c_{n,r,j}(f)&=\frac{(-1)^{r+j}r!}{j!}
   [t^{r-j}]\left(\frac{f(t)}t\right)^{-r-1}
       \alpha_{n,r}\!\left(\log\frac{f(t)}t\right).
       \label{nc:coefficient}
\end{align}
The coordinate jet through degree $r+1$ suffices for every Stieltjes
index $n$.
\end{theorem}

\begin{proof}
Differentiate $x^{z-1}$ in its argument:
\[
 \partial_x^rx^{z-1}
    =(-1)^rr!\mathcal E_r(z)x^{z-r-1}.
\]
Taking $n![z^n]$ gives the exact singular polynomial
\begin{equation}\label{nc:singular-polynomial}
 \partial_x^r\frac{\log^n x}{x}
  =(-1)^rr!x^{-r-1}
      \sum_{j=0}^{\min(r,n)}
       \frac{(-1)^je_j(r)n!}{(n-j)!}\log^{n-j}x.
\end{equation}
Apply Theorem~\ref{nc:universal} term by term. Its factor
$1/(n-j+1)$ converts the displayed polynomial into
\eqref{nc:alpha-finite}, while the analytic remainder in
\eqref{nc:Stieltjes-local} contributes no residue. Expanding the test
function through degree $r$ yields \eqref{nc:coefficient}, since
$\langle\delta_0^{(j)},\phi\rangle=(-1)^j\phi^{(j)}(0)$.
The coefficient in that formula uses the expansion of $f(t)/t$ only
through degree $r$, proving the last assertion.
\end{proof}

For comparison, distributional differentiation gives
\begin{equation}\label{nc:derivative-contact}
 D^rG_n=T_{n,r}+b_{n,r}(0)\delta_0^{(r)},\qquad
 b_{n,r}(0)=(-1)^{n+1}n!e_{n+1}(r).
\end{equation}
This also follows by the regulator computation: at identity scale the
polar numerator for $\partial_x^rx^{z-1}$ is
$\mathcal E_r(z)\delta_0^{(r)}$. Thus
\begin{align}
 f^*D^rG_n&=\FP_t\gamma_n^{(r)}(f(t))+\mathfrak C_{n,r}(f),
       \label{nc:distributional-law}\\
 \langle\mathfrak C_{n,r}(f),\phi\rangle
   &=(-1)^rr!\Res_{t=0}
       f(t)^{-r-1}b_{n,r}(\lambda_f(t))\phi(t)\dd t.
       \label{nc:Cresidue}
\end{align}
The finite coefficient formula for $\mathfrak C$ is
\eqref{nc:coefficient} with $\alpha$ replaced by $b$.

\begin{corollary}[Recovery of the affine contact polynomial]
\label{nc:affine}
At $f(t)=qt$, $q>0$,
\begin{align}
 \mathfrak A_{n,r}(qt)&=q^{-r-1}\alpha_{n,r}(\log q)
                         \delta_0^{(r)},\label{nc:affine-A}\\
 D_t^r f^*G_n
   &=\FP_t\bigl(q^r\gamma_n^{(r)}(qt)\bigr)
        +\frac{b_{n,r}(\log q)}q\delta_0^{(r)}.
       \label{nc:affine-C}
\end{align}
The second identity is exactly the local form of the incoming
integer-dilation theorem, including its normalization by $1/q$.
\end{corollary}
\begin{proof}
The function $f(t)/t$ is constant. Only the $j=r$ term survives in
\eqref{nc:coefficient}; also
$D_t^rf^*G_n=q^rf^*D^rG_n$.
\end{proof}

\subsection{Explicit corrections through the second derivative}

Write
\begin{equation}\label{nc:f-expansion}
 f(t)=qt(1+a t+b t^2+O(t^3)),\qquad L=\log q.
\end{equation}
Thus $a=f''(0)/(2q)$ and $b=f'''(0)/(6q)$. The first coefficient
polynomials are
\begin{align}
 \alpha_{0,r}(L)&=L,\notag\\
 \alpha_{1,r}(L)&=\frac{L^2}{2}-H_rL,\label{nc:first-alpha}\\
 \alpha_{2,r}(L)&=\frac{L^3}{3}-H_rL^2
                    +(H_r^2-H_r^{(2)})L.\notag
\end{align}
Here $H_0=H_0^{(2)}=0$. At $r=0$, the entire answer is
\begin{equation}\label{nc:undifferentiated}
 \mathfrak A_{n,0}(f)=\frac{L^{n+1}}{q(n+1)}\delta_0.
\end{equation}
Nonlinear coefficients first affect argument derivatives.

For $r=1$ set $\alpha=\alpha_{n,1}(L)$; primes below differentiate
this polynomial in $L$. Equation~\eqref{nc:coefficient} becomes
\begin{equation}\label{nc:r1}
 \mathfrak A_{n,1}(f)
  =q^{-2}\left[\alpha\delta_0'
                   +a(2\alpha-\alpha')\delta_0\right].
\end{equation}
For $r=2$, with $\alpha=\alpha_{n,2}(L)$, it gives
\begin{align}\label{nc:r2}
 \mathfrak A_{n,2}(f)=q^{-3}\bigl[&\alpha\delta_0''
       +2a(3\alpha-\alpha')\delta_0'\notag\\
       &+\{(12a^2-6b)\alpha+(2b-7a^2)\alpha'
                                      +a^2\alpha''\}\delta_0\bigr].
\end{align}
To verify these formulas directly, use
\[
 \lambda_f=L+a t+(b-a^2/2)t^2+O(t^3),\qquad
 (1+a t+b t^2)^{-3}
      =1-3a t+(6a^2-3b)t^2+O(t^3).
\]
The same displays with $b_{n,r}$ in place of $\alpha_{n,r}$ give the
correction $\mathfrak C_{n,r}$ for distributional derivatives.

At unit tangent, these are particularly short:
\begin{equation}\label{nc:r1unit}
 \mathfrak A_{n,1}=
 \begin{cases}
  -a\delta_0,&n=0,\\
  a\delta_0,&n=1,\\
  0,&n\geq2,
 \end{cases}
\end{equation}
and
\begin{equation}\label{nc:r2unit}
 \mathfrak A_{n,2}=
 \begin{cases}
  -2a\delta_0'+(2b-7a^2)\delta_0,&n=0,\\
  3a\delta_0'+(-3b+\tfrac{23}{2}a^2)\delta_0,&n=1,\\
  -2a\delta_0'+(2b-10a^2)\delta_0,&n=2,\\
  3a^2\delta_0,&n=3,\\
  0,&n\geq4.
 \end{cases}
\end{equation}
These expressions retain curvature and the cubic coordinate coefficient
even though the linear scale has disappeared.

\begin{example}[A direct cutoff check of the curvature terms]
\label{nc:direct-cutoff}
Take $f(t)=t+a t^2+b t^3+O(t^4)$ and $\gamma_0''(x)$, whose singular
part is $2x^{-3}$. Its inverse is
$g(x)=x-a x^2+(2a^2-b)x^3+O(x^4)$. For a test function equal to one
near zero, pullback uses the inverse Jacobian
$g'(x)=1-2a x+(6a^2-3b)x^2+O(x^3)$.
A primitive of the terms that can contribute a finite constant is
\[
 -x^{-2}+4a x^{-1}+(12a^2-6b)\log x.
\]
It has zero constant in the $x$ coordinate. After $x=f(\eps)$ its
constant is $(2b-3a^2)-4a^2=2b-7a^2$, giving the $\delta_0$
coefficient in \eqref{nc:r2unit}. For a test function equal to $t$,
the inverse multiplier is $g(x)g'(x)=x-3a x^2+O(x^3)$; the relevant
primitive is $-2/x-6a\log x$. Its composed constant is $2a$, which
equals the action of $-2a\delta_0'$ on that test function.
This calculation uses cutoff primitives directly and confirms both
nonlinear signs without invoking the residue formula.
\end{example}

\subsection{Sharp invariance and the order of the local correction}

\begin{theorem}[Sharp Stieltjes-index filtration]
\label{nc:rigidity}
Suppose $f'(0)=1$, put $a=f''(0)/2$, and let $r\geq1$. Then
$\mathfrak A_{n,r}(f)$ is a linear combination of
$\delta_0^{(j)}$ with
\begin{equation}\label{nc:order-filtration}
 0\leq j\leq\min\{r-1,\,2r-n-1\}.
\end{equation}
An empty range means the zero distribution. In particular,
\begin{equation}\label{nc:high-invariance}
 \boxed{\quad
 f^*\FP_x\gamma_n^{(r)}
      =\FP_t\gamma_n^{(r)}(f(t))\qquad(n\geq2r).
 \quad}
\end{equation}
For $r\leq n\leq2r-1$, let
$d=n+1-r$ and $k=2r-n-1$. The coefficient of the highest possible
derivative is
\begin{equation}\label{nc:sharp-coefficient}
 [\delta_0^{(k)}]\mathfrak A_{n,r}(f)
        =\frac{(-1)^kn!}{d!\,k!}a^d.
\end{equation}
For $0\leq n\leq r$, the coefficient of $\delta_0^{(r-1)}$ is
\begin{equation}\label{nc:lower-sharp}
 [\delta_0^{(r-1)}]\mathfrak A_{n,r}(f)
       =(-1)^{n+1}r\,n!e_n(r)a.
\end{equation}
The two formulas agree at $n=r$. If $a\ne0$, every stated upper
order is attained. At the last nonzero index,
\begin{equation}\label{nc:last-index}
 \mathfrak A_{2r-1,r}(f)
     =\frac{(2r-1)!}{r!}a^r\delta_0.
\end{equation}
For $r=0$, every correction vanishes at unit tangent by
\eqref{nc:undifferentiated}.
\end{theorem}

\begin{proof}
Now $\lambda_f(t)=a t+O(t^2)$. Formula~\eqref{nc:alpha-finite}
shows that the smallest possible power of $L$ in $\alpha_{n,r}$ is
$\max(1,n+1-r)$. Therefore the analytic factor multiplying
$t^{-r-1}$ in \eqref{nc:Stieltjes-residue} vanishes to at least that
order. Its residue against $\phi$ can use only derivatives of orders
at most $r-\max(1,n+1-r)$, proving
\eqref{nc:order-filtration} and \eqref{nc:high-invariance}.

For $r\leq n\leq2r-1$ the lowest term is, using $e_r(r)=1/r!$,
\[
 \alpha_{n,r}(L)
    =\frac{(-1)^rn!}{r!\,d!}L^d+O(L^{d+1}).
\]
To produce $\delta_0^{(k)}$ in \eqref{nc:coefficient}, only the
coefficient of $t^{r-k}=t^d$ is required. It is the displayed
coefficient times $a^d$; the constant term of $(f(t)/t)^{-r-1}$ is
one. Substitution proves \eqref{nc:sharp-coefficient}. For $n\leq r$
the linear coefficient of $\alpha_{n,r}$ is
$(-1)^nn!e_n(r)$. Taking $j=r-1$ in the same formula proves
\eqref{nc:lower-sharp}. All these elementary symmetric coefficients
are nonzero in their indicated ranges. Equation~\eqref{nc:last-index}
is the case $d=r,k=0$.
\end{proof}

\begin{remark}[Higher tangency gives stronger invariance]
\label{nc:higher-tangency}
If $f(t)=t+O(t^{d+1})$, then $\lambda_f=O(t^d)$. If $d>r$, every
$\mathfrak A_{n,r}$ vanishes. If $1\leq d\leq r$, it suffices that
\[
 n\geq r+\lfloor r/d\rfloor.
\]
This threshold is sharp in general. For
$f(t)=t+c t^{d+1}+O(t^{d+2})$, set
$h=\lfloor r/d\rfloor$, $n=r+h-1$ and $k=r-dh$. The coefficient of
$\delta_0^{(k)}$ is
$(-1)^kn!c^h/(h!k!)$. This follows from the same lowest-power argument.
\end{remark}

\subsection{Polygamma functions, full derivatives, and pointwise products}

For polygamma, the correction in
\eqref{nc:Stieltjes-law} is $-\mathfrak A_{0,r}$ because
$\gamma_0^{(r)}=-\psi^{(r)}$. For example, at
$f(t)=t+a t^2+O(t^3)$,
\begin{equation}\label{nc:trigamma}
 f^*\FP_x\psi'(x)
    =\FP_t\psi'(f(t))+a\delta_0.
\end{equation}
The undifferentiated function $\log\Gamma(x)$ is locally integrable;
its logarithmic singularity has no negative-power coefficient in
\eqref{nc:loggerm}. Its ordinary pullback therefore has no anomaly.

A derivative in the argument of $\gamma_n$ differs from a full
derivative of $\gamma_n(f(t))$. The latter is also finite and explicit.
Let $B_{k,r}$ denote the partial exponential Bell polynomial in the
ordinary Fa\`a di Bruno formula. For $k\geq1$,
\begin{align}\label{nc:chain-rule}
 D_t^kf^*G_n-&\FP_tD_t^k\bigl(\gamma_n(f(t))\bigr)\notag\\
   &=\sum_{r=1}^k
       B_{k,r}(f',f'',\ldots,f^{(k-r+1)})\,
                         \mathfrak C_{n,r}(f).
\end{align}
Indeed, the distributional chain rule holds first for smooth functions
and then by pullback and differentiation of distributions; it gives
the displayed Bell sum with $f^*D^rG_n$. Finite parts commute with
multiplication by a smooth function, directly from their test-function
definition, so \eqref{nc:distributional-law} proves
\eqref{nc:chain-rule}. To put the answer in constant delta coordinates,
use the finite identity
\begin{equation}\label{nc:delta-multiplier}
 v\delta_0^{(j)}
   =\sum_{\ell=0}^j(-1)^{j-\ell}\binom j\ell
       v^{(j-\ell)}(0)\delta_0^{(\ell)}.
\end{equation}
Thus no indefinite distribution operation remains in the formula.

The universal theorem also treats a pointwise product once its own
finite part has been specified. This gives a small exact illustration
relevant to collision normalizations.

\begin{corollary}[Coordinate law for a coincident Stieltjes product]
\label{nc:coincident}
Let $h(x)=\gamma_m(x)\gamma_n(x)$, $M=m+n$, and use
\eqref{nc:f-expansion}. Then
\begin{align}\label{nc:product-anomaly}
 \mathfrak A_f[h]={}&-\frac{L^{M+1}}{q^2(M+1)}\delta_0'\notag\\
 &+\biggl[
  \frac{a}{q^2}\left(L^M-\frac{2L^{M+1}}{M+1}\right)
  +\frac{\gamma_n(1)L^{m+1}}{q(m+1)}
  +\frac{\gamma_m(1)L^{n+1}}{q(n+1)}
                  \biggr]\delta_0.
\end{align}
As usual $L^0=1$, also when $L=0$. At unit tangent the sole nonzero
case is $m=n=0$, with correction $a\delta_0$.
\end{corollary}
\begin{proof}
Equation~\eqref{nc:Stieltjes-local} gives the complete nonintegrable
part
\[
 \frac{\log^{M}x}{x^2}
 +\frac{\gamma_n(1+x)\log^m x+
         \gamma_m(1+x)\log^n x}{x}.
\]
For the first term, expand $f^{-2}\lambda_f^{M+1}/(M+1)$ through
the coefficient of $t^{-1}$. It gives the first line of
\eqref{nc:product-anomaly} and its $a/q^2$ term. The two simple-pole
terms require only the values $\gamma_n(1)$ and $\gamma_m(1)$,
and yield the remaining terms. The analytic product remainder has no
residue.
\end{proof}

This corollary compares coordinate choices for a specified finite part
of an ordinary pointwise product. It does not identify that finite
part with a collision limit of separated products, and it does not
define multiplication of arbitrary distributions at common singular
supports. It supplies explicit local normalization data for such a
future comparison.

\subsection{Exact algorithm and independent checks}

For fixed $n,r$, the complete computation consists of four finite
operations: form $\mathcal E_r$, extract $\alpha_{n,r}$, truncate
$\log(f(t)/t)$ and $(f(t)/t)^{-r-1}$ through degree $r$, and read the
coefficients in \eqref{nc:coefficient}. The coefficients are rational
polynomials in the normalized coordinate derivatives and $\log q$,
with the overall factor $q^{-r-1}$. Composition is checked by the
binomial identity in Theorem~\ref{nc:composition}.

The accompanying program \src{code/verify_nonlinear.py} performs exact
rational symbolic checks of the coefficient generator, the scale-aware
cocycle, the generic first- and second-derivative formulas, and the sharp
filtration. It also compares the correction with an independent direct
cutoff-primitive calculation for finite test jets. The record is
\src{results/nonlinear_checks.json}. These finite tests check the
implementation and normalization; the all-index analytic assertions
are established by the proofs above.
